\documentclass[10pt,a4paper]{amsart}
\usepackage{amsmath,amssymb,amsthm,graphicx}
\usepackage{verbatim}
\newenvironment{reference}{\comment}{\endcomment}
%\newenvironment{reference}{}{}
\newenvironment{slogan}{\comment}{\endcomment}
\newenvironment{history}{\comment}{\endcomment}

% For commutative diagrams we use Xy-pic
\usepackage[all]{xy}

% We use 2cell for 2-commutative diagrams.
\xyoption{2cell}
% Theorem environments.
%
\theoremstyle{plain}
\newtheorem{theorem}[subsection]{Theorem}
\newtheorem{proposition}[subsection]{Proposition}
\newtheorem{lemma}[subsection]{Lemma}

\theoremstyle{definition}
\newtheorem{definition}[subsection]{Definition}
\newtheorem{example}[subsection]{Example}
\newtheorem{exercise}[subsection]{Exercise}
\newtheorem{situation}[subsection]{Situation}

\theoremstyle{remark}
\newtheorem{remark}[subsection]{Remark}
\newtheorem{remarks}[subsection]{Remarks}

\numberwithin{equation}{subsection}

% Macros
%
\def\lim{\mathop{\mathrm{lim}}\nolimits}
\def\colim{\mathop{\mathrm{colim}}\nolimits}
\def\Spec{\mathop{\mathrm{Spec}}}
\def\Hom{\mathop{\mathrm{Hom}}\nolimits}
\def\Ext{\mathop{\mathrm{Ext}}\nolimits}
\def\SheafHom{\mathop{\mathcal{H}\!\mathit{om}}\nolimits}
\def\SheafExt{\mathop{\mathcal{E}\!\mathit{xt}}\nolimits}
\def\Sch{\mathit{Sch}}
\def\Mor{\mathop{\mathrm{Mor}}\nolimits}
\def\Ob{\mathop{\mathrm{Ob}}\nolimits}
\def\Sh{\mathop{\mathit{Sh}}\nolimits}
\def\NL{\mathop{N\!L}\nolimits}
\def\CH{\mathop{\mathrm{CH}}\nolimits}
\def\proetale{{pro\text{-}\acute{e}tale}}
\def\etale{{\acute{e}tale}}
\def\QCoh{\mathit{QCoh}}
\def\Ker{\mathop{\mathrm{Ker}}}
\def\Im{\mathop{\mathrm{Im}}}
\def\Coker{\mathop{\mathrm{Coker}}}
\def\Coim{\mathop{\mathrm{Coim}}}

% Boxtimes
%
\DeclareMathSymbol{\boxtimes}{\mathbin}{AMSa}{"02}

%
% Macros for moduli stacks/spaces
%
\def\QCohstack{\mathcal{QC}\!\mathit{oh}}
\def\Cohstack{\mathcal{C}\!\mathit{oh}}
\def\Spacesstack{\mathcal{S}\!\mathit{paces}}
\def\Quotfunctor{\mathrm{Quot}}
\def\Hilbfunctor{\mathrm{Hilb}}
\def\Curvesstack{\mathcal{C}\!\mathit{urves}}
\def\Polarizedstack{\mathcal{P}\!\mathit{olarized}}
\def\Complexesstack{\mathcal{C}\!\mathit{omplexes}}
% \Pic is the operator that assigns to X its picard group, usage \Pic(X)
% \Picardstack_{X/B} denotes the Picard stack of X over B
% \Picardfunctor_{X/B} denotes the Picard functor of X over B
\def\Pic{\mathop{\mathrm{Pic}}\nolimits}
\def\Picardstack{\mathcal{P}\!\mathit{ic}}
\def\Picardfunctor{\mathrm{Pic}}
\def\Deformationcategory{\mathcal{D}\!\mathit{ef}}
\usepackage{cleveref}
\crefname{originallemma}{original lemma}{original lemmas}
\Crefname{originallemma}{Original lemma}{Original lemmas}
\makeatletter
\newlabel{lemma-characterize-strict-general}{{05SI}{source}{}{sourcepointer-1}{}}
\newlabel{lemma-characterize-strict-general@cref}{{[originallemma][0][]05SI}{[1][0][]source}{}{}{}}
\makeatother
\crefname{originallemma}{original lemma}{original lemmas}
\Crefname{originallemma}{Original lemma}{Original lemmas}
\makeatletter
\newlabel{lemma-ses-gr}{{05SP}{source}{}{sourcepointer-2}{}}
\newlabel{lemma-ses-gr@cref}{{[originallemma][0][]05SP}{[1][0][]source}{}{}{}}
\makeatother
\title{Characterizations of strict morphisms of finite filtered objects}
\author{The Stacks Project Authors}
\date{stacks-project:homology.tex:local-snapshot-20261012; GNU FDL 1.2 or later}
\begin{document}
\maketitle
\noindent\textit{Editorial note.} All delivered mathematical text is the author's own, copied from the pinned Stacks Project chapter file homology.tex; no statement or proof step was rewritten, reordered or removed.\par\medskip
\noindent\textit{Editorial note.} The proof is the author's complete argument for all six equivalences: two of them are attributed to named same-chapter lemmas carried on this card as navigation stubs, and the one substantive implication is written out in full inside the counted span with its descending induction and five-lemma step.\par\medskip
\noindent\textit{Editorial note.} Every labelled reference occurring inside the counted spans is carried in assembly.external\_references as a non-load-bearing navigation pointer (same-chapter ones at their definition site, cross-chapter ones at the invocation inside the counted span), so no load-bearing mathematics is replaced by a stub. No bibliography citation occurs inside the counted spans.\par\medskip
\noindent\textit{Editorial note.} The same-name check for this card scanned every staged Stacks chapter of this repository -- all sources/stacks-*/pinned-original-source.tex files, the card's own chapter excluded -- and the labels found are recorded in dedup.same\_name\_cross\_chapter on the card.\par\medskip
\noindent\textit{Editorial note.} The preamble is the chapter's own native preamble excerpt already used for this source: the comment-type reference, slogan and history environments (native preamble lines 8-13), the Xy-pic setup (native preamble lines 14-18) and the native theorem declarations with the mathematical macros (native preamble lines 154-217, the actual last line of the file).\par\medskip
\noindent\textit{Allowed standard prerequisites (editorial).}\begin{itemize}
\item The general characterization of strictness. Same chapter, Lemma lemma-characterize-strict-general. Supplies the equivalence between strictness of a morphism and the canonical map from its coimage to its image being an isomorphism.
\item Short exact sequences and associated graded objects. Same chapter, Lemma lemma-ses-gr. Supplies that the exactness of the sequences of graded pieces is equivalent to the graded pieces of coimage and image being isomorphic, which gives the equivalence of (3) with (4), (5) and (6).
\item The five lemma. Same chapter, the five lemma for abelian categories. Supplies the conclusion of the descending induction from the commutative diagram whose rows are short exact sequences.
\end{itemize}\medskip
\ifdefined\hypertarget\hypertarget{sourcepointer-1}{}\fi
\noindent\textit{Source reference.} Original lemma 05SI: sources/stacks-homology/pinned-original-source.tex, source lines 4969--4969. Navigation pointer invoked at proof line 5263; the counted span contains the pointer only and the cited result is not reproduced inside it, so the pointer is registered as a non-load-bearing navigation stub in the established form for Stacks material. No bibliography citation is involved.\par\medskip
\ifdefined\hypertarget\hypertarget{sourcepointer-2}{}\fi
\noindent\textit{Source reference.} Original lemma 05SP: sources/stacks-homology/pinned-original-source.tex, source lines 5199--5199. Navigation pointer invoked at proof line 5265; the counted span contains the pointer only and the cited result is not reproduced inside it, so the pointer is registered as a non-load-bearing navigation stub in the established form for Stacks material. No bibliography citation is involved.\par\medskip
\begin{lemma}
\label{lemma-characterize-strict}
Let $\mathcal{A}$ be an abelian category.
Let $f : A \to B$ be a morphism of finite
filtered objects of $\mathcal{A}$. The following are equivalent
\begin{enumerate}
\item $f$ is strict,
\item the morphism $\Coim(f) \to \Im(f)$ is an isomorphism,
\item $\text{gr}(\Coim(f)) \to \text{gr}(\Im(f))$ is an
isomorphism,
\item the sequence
$\text{gr}(\Ker(f)) \to \text{gr}(A) \to \text{gr}(B)$
is exact,
\item the sequence $\text{gr}(A) \to \text{gr}(B) \to
\text{gr}(\Coker(f))$ is exact, and
\item the sequence
$$
0 \to
\text{gr}(\Ker(f)) \to
\text{gr}(A) \to
\text{gr}(B) \to
\text{gr}(\Coker(f)) \to 0
$$
is exact.
\end{enumerate}
\end{lemma}

\begin{proof}
The equivalence of (1) and (2) is
Lemma \ref{lemma-characterize-strict-general}.
By
Lemma \ref{lemma-ses-gr}
we see that (4), (5), (6) imply (3) and that (3) implies (4), (5), (6).
Hence it suffices to show that (3) implies (2).
Thus we have to show that if $f : A \to B$ is an injective and surjective
map of finite filtered objects which induces and isomorphism
$\text{gr}(A) \to \text{gr}(B)$, then $f$ induces an isomorphism of
filtered objects. In other words, we have to show that
$f(F^pA) = F^pB$ for all $p$.
As the filtrations are finite we may prove this by descending induction
on $p$. Suppose that $f(F^{p + 1}A) = F^{p + 1}B$.
Then commutative diagram
$$
\xymatrix{
0 \ar[r] &
F^{p + 1}A \ar[r] \ar[d]^f &
F^pA \ar[r] \ar[d]^f &
\text{gr}^p(A) \ar[r] \ar[d]^{\text{gr}^p(f)} &
0 \\
0 \ar[r] &
F^{p + 1}B \ar[r] &
F^pB \ar[r] &
\text{gr}^p(B) \ar[r] &
0
}
$$
and the five lemma imply that $f(F^pA) = F^pB$.
\end{proof}

\end{document}
